% Short report -- Time Series (SS 2026), Deliverable 3
% Uses the ICASSP-2026 style provided with the course template.
\documentclass{article}
\usepackage{spconf,amsmath,graphicx,hyperref,booktabs,multirow}

\title{Parkinson's Disease Detection from Gait Force Signals\\ with a State-Space (Mamba) Model:\\ Results and the Role of the Evaluation Protocol}

\name{Ahmadreza Nourozi\thanks{Ahmadreza Nourozi, ahmadrezanourozii@gmail.com,
matriculation number 23726011, M.Sc.\ Artificial Intelligence.
\emph{Please complete: home address, date of birth.}}}
\address{Pattern Recognition Lab, Friedrich\mbox{-}Alexander\mbox{-}Universit\"at Erlangen\mbox{-}N\"urnberg, Erlangen, Germany}

\begin{document}
\ninept
\maketitle

\begin{abstract}
Vertical ground reaction force (VGRF) recordings offer an inexpensive route to
objective Parkinson's disease (PD) screening, but reported accuracies vary
enormously because evaluation protocols differ. We classify PD versus healthy
control (HC) from the PhysioNet gait database (100 subjects, 16 force sensors)
using a Mamba state-space model, and we evaluate it under two protocols on
identical data. Under a strictly subject-wise protocol, in which all windows of
a subject stay inside one split and the configuration is selected on validation
rather than on test, the model reaches $75 \pm 4$\,\% accuracy, $0.75$ weighted
F1 and $0.85$ AUC on both feet, slightly ahead of the Random Forest, SVM-RBF,
1D-ResNet and LSTM baselines. An ablation shows the gain comes from describing
each second by its channel statistics rather than feeding raw samples, which
also removes a training collapse observed on two folds, while a bidirectional
variant of the causal scan changes nothing. Re-evaluating the same features and baselines
under the window-wise protocol that is common in the literature, where windows
of one subject appear in training and test, raises Random Forest accuracy from
$73$\,\% to $98$\,\% and SVM accuracy from $68$\,\% to $88$\,\% -- an
inflation of twenty to twenty-five points caused solely by the split. We
report the subject-wise numbers as our result and the window-wise numbers only
as a cautionary reference.
\end{abstract}

\begin{keywords}
Parkinson's disease, gait analysis, state-space models, Mamba, data leakage
\end{keywords}

\section{Introduction}
\label{sec:intro}

Parkinson's disease (PD) is the second most common neurodegenerative disorder,
and disturbed gait is one of its earliest and most visible motor signs. Clinical
scales such as the UPDRS rely on expert observation and are coarsely quantised,
which is why sensor-based measures are attractive. Instrumented insoles record
the vertical ground reaction force (VGRF) under each foot cheaply and at high
rate, and the resulting series carries the stride-timing irregularity,
shortened swing phase and left--right asymmetry that characterise parkinsonian
gait \cite{hausdorff2007, frenkel2005treadmill}.

Accuracies published on the PhysioNet gait database span roughly 70--99\,\%,
much wider than the differences one would expect between the model families
involved. A two-minute recording yields more than a hundred overlapping
windows, and if those are divided at random then windows of one subject end up
in both training and test. A classifier can then succeed by recognising the
individual instead of the disease -- a textbook case of leakage
\cite{kaufman2012leakage}, and a failure mode repeatedly documented in medical
machine learning \cite{varoquaux2022machine}. This report therefore asks two
questions at once: how well a state-space model classifies these signals, and
how much of a reported accuracy is an artefact of the split.

\subsection{Related work}
\label{ssec:sota}
Classical pipelines on this database pair hand-crafted statistics of the VGRF
channels with Random Forests or kernel SVMs. Deep models -- 1D-ResNets
\cite{he2016resnet}, recurrent networks and more recently transformers -- take
either the raw signal or sequences of feature vectors. Structured state-space
models are a newer alternative: S4 \cite{gu2022s4} and then Mamba
\cite{gu2023mamba} reach transformer-level sequence modelling at linear rather
than quadratic cost in sequence length, because Mamba makes the state transition
depend on the input instead of being fixed. The architecture has been applied to
speech, EEG and ECG; we are not aware of a VGRF-based PD study that uses it
under a subject-disjoint protocol.

\subsection{Contributions of this work}
\label{ssec:contr}
We train a Mamba classifier on this dataset and compare it with the Random
Forest, SVM-RBF, 1D-ResNet and LSTM baselines from the previous stage of the
project, evaluating every model with the same code and the same subject-disjoint
folds. Two findings came out of the experiments rather than out of the plan. The
input representation mattered more than the architecture: feeding raw samples
made training collapse on two of five folds, and describing each second by its
channel statistics both removed the collapse and improved accuracy. And the
split unit mattered more than either: re-running the identical pipeline with
window-wise splits raises Random Forest accuracy from 73\,\% to 98\,\%, which we
report as a caution rather than as a result.

\section{Materials \& Methods}
\label{sec:matmeth}

Figure~\ref{fig:pipeline} summarises the complete pipeline, from the raw force
recordings to the per-subject PD/HC decision.

\begin{figure*}[t]
  \centering
  \includegraphics[width=\textwidth]{figures/fig1_pipeline}
  \caption{End-to-end pipeline. Raw VGRF recordings (16 sensors plus two
  per-foot totals, 100\,Hz) are trimmed, windowed and converted into token
  sequences; a bidirectional Mamba encoder produces a window score, and the
  scores of a subject's windows are averaged into one decision. Normalisation
  statistics and every hyper-parameter are estimated on training and
  validation data only, and each of the five subject-wise folds contributes a
  disjoint set of test subjects.}
  \label{fig:pipeline}
\end{figure*}

\subsection{Data}
\label{ssec:data}
We use the PhysioNet \emph{Gait in Parkinson's Disease} database, pooling its
three studies (Ga, Ju, Si). Each subject walked on level ground for about two
minutes while eight force sensors per foot sampled at 100\,Hz; the dataset
provides the 16 sensor channels plus the total force under each foot. Our
subset comprises the 100 subjects with a normal-walk recording, balanced by
design into 50 PD patients and 50 healthy controls. The course-provided
five-fold partition assigns 60 training, 20 validation and 20 test subjects
per fold, class-balanced, with test sets disjoint across folds so that pooling
the five test folds covers every subject exactly once.

\subsubsection{Pre-processing}
The first five seconds of each recording are discarded to remove gait
initiation. The remainder is cut into 5\,s windows with a 1\,s step
(500 samples per window), yielding 10\,407 windows in total, between 35 and
254 per subject. Each window inherits its subject's label. Channel
normalisation is a per-channel z-score whose mean and standard deviation are
computed over the training windows of the current fold only and then applied
unchanged to validation and test windows.

Three token representations were compared, all sharing the above windowing:
\textbf{(i) raw samples}, each window a $18 \times 500$ matrix;
\textbf{(ii) amplitude-statistic tokens}, one token per second holding 11
statistics (mean, standard deviation, skewness, kurtosis, RMS, zero-crossing
rate, median, minimum, maximum, spectral energy, spectral entropy) for each of
the 18 channels, i.e.\ 198 dimensions per token; and
\textbf{(iii) stride-cycle tokens}, one token per gait cycle carrying stride,
stance and swing time, stance ratio, normalised peak force and impulse,
stride-time increment and double-support time for both feet. Foot contacts are
detected by thresholding the per-foot total force at 5\,\% of its maximum, and
strides outside 0.5--2.5\,s are rejected as implausible; the 100 subjects give
between 34 and 166 strides each.

\subsubsection{Materials/Instruments used}
Models were trained on a Kaggle NVIDIA Tesla T4 (16\,GB, compute capability
7.5) with PyTorch 2.10 (CUDA 12.8) and \texttt{mamba-ssm} 2.3.2, using mixed
precision. Development, feature extraction and evaluation ran locally on an
Apple M-series machine. Code was version-controlled on GitHub and pulled by
the Kaggle notebook at the start of every run, so each reported number comes
from a known commit.

\subsection{Proposed approach}
\label{ssec:mod}

\subsubsection{Bidirectional Mamba classifier}
A Mamba block is a selective state-space layer: its state transition and input
projections are functions of the input, so the recurrence can forget or retain
information conditioned on content, and it runs in linear time in the sequence
length \cite{gu2023mamba}. Standard Mamba is causal, which suits language
modelling but wastes information when an entire window is available at
inference. We therefore use a bidirectional block: the normalised token
sequence is processed by a forward mixer and, after reversal, by a second
mixer with independent parameters, and the two outputs are summed before the
residual connection.

The classifier projects tokens to $d_{\text{model}}$, applies $L$ such blocks,
and pools over time by concatenating the mean and the maximum of the final
representation, which a linear layer maps to two logits. Three sizes were
studied ($d_{\text{model}}, L$) $\in \{(64,2), (128,4), (256,6)\}$; the
reported model is the largest.

\subsubsection{Training strategy}
Per fold we grid-search dropout $\in \{0.1, 0.3\}$ and learning rate
$\in \{3\!\cdot\!10^{-4}, 10^{-4}\}$, training with AdamW (weight decay
$10^{-3}$), three warm-up epochs followed by cosine decay, batch size 128 and
at most 30 epochs. The configuration with the highest validation ROC-AUC is
retrained on the union of training and validation subjects for the epoch count
that was optimal on validation, and that model predicts the test subjects --
the deep-learning analogue of \texttt{GridSearchCV(refit=True)} on a
predefined split. Early stopping uses validation ROC-AUC with patience 6 and a
minimum of 8 epochs. Because some folds occasionally converged to a degenerate
solution, a fold whose best validation AUC stays below 0.70 is retrained with
a different seed (at most three attempts). The final model averages the window
scores of three seeds.

\subsubsection{Decision rule and subject-level aggregation}
The network scores windows, but the clinical unit is the subject. A subject's
score is the mean of its window scores. The decision threshold is calibrated
per fold on that fold's \emph{validation} subjects by maximising balanced
accuracy, using the model trained on training data only; test data never
influences it. Results at the fixed threshold 0.5 are also reported and differ
by at most two points.

\subsubsection{Evaluation metrics}
We report accuracy, precision, recall and F1 (weighted by class support, so
that recall equals accuracy by construction), together with sensitivity,
specificity and ROC-AUC, plus confusion matrices. All values are the mean and
standard deviation over 1000 bootstrap resamples of the pooled test
predictions -- resampling subjects for subject-level metrics and windows for
window-level metrics.

\section{Experiments \& Results}
\label{sec:exp}

\textbf{Partitioning.} Unless stated otherwise, every experiment uses the
five predefined subject-wise folds (60/20/20 subjects). Training uses the 60
training subjects, hyper-parameters and the decision threshold are selected on
the 20 validation subjects, and the 20 test subjects are predicted once. The
five test sets are disjoint, so pooling them yields exactly one prediction per
subject and 10\,407 window predictions in total. This is cross-validation in
the sense that every subject is tested exactly once, without any subject
appearing in two roles.

\textbf{Experiments performed.} (1) token representation: raw samples versus
amplitude statistics versus stride cycles versus their fusion; (2) architecture:
unidirectional versus bidirectional, three encoder sizes, regularisation and
seed ensembling; (3) foot configuration: left, right, both; (4) protocol
comparison: subject-wise versus window-wise splits. Fifteen GPU runs were
executed in total.

\subsection{Experiment 1: token representation}
Table~\ref{tab:tokens} reports subject-level accuracy on both feet for the
input representations. Raw signals were markedly unstable: on folds 1 and 4
the validation AUC stayed near chance ($0.53$ and $0.54$) while folds 2, 3 and
5 exceeded $0.92$, and neither longer windows (30\,s, 60\,s), seed restarts
nor a larger encoder repaired it. Amplitude-statistic tokens removed the
instability -- the weakest fold rose to $0.75$ validation AUC -- and improved
accuracy by seven points. Stride-cycle tokens alone were weaker, and fusing
them with amplitude tokens did not surpass amplitude tokens alone, although at
subject level the two families were complementary in a Random Forest analysis
(AUC $0.79$ stride only, $0.83$ amplitude only, $0.86$ combined).

\begin{table}[t]
  \centering
  \caption{Effect of the input representation on subject-level accuracy (both
  feet, subject-wise protocol). ``Worst fold'' is the lowest validation ROC-AUC
  over the five folds and exposes the training collapse seen with raw input.}
  \label{tab:tokens}
  \begin{tabular}{lcc}
    \toprule
    Input representation & Accuracy (\%) & Worst fold AUC \\
    \midrule
    Raw samples, 5\,s windows            & 69 & 0.53 \\
    Raw samples, 30\,s windows           & 69 & 0.63 \\
    Raw samples, 60\,s windows           & 65 & 0.59 \\
    Stride-cycle tokens                  & 66 & 0.71 \\
    Amplitude ++ stride (fused)          & 78 & 0.71 \\
    \textbf{Amplitude statistics}        & \textbf{80} & \textbf{0.77} \\
    \bottomrule
  \end{tabular}
\end{table}

\subsection{Experiment 2: architecture ablation}
With amplitude tokens fixed, encoder size was the decisive architectural
factor: enlarging the encoder from $d_{\text{model}}=128$, $L=4$ to
$d_{\text{model}}=256$, $L=6$ raised subject-level accuracy from 72\,\% to
80\,\%.

The bidirectional extension did \emph{not} reproduce that gain. At
$d_{\text{model}}=128$ the causal and the bidirectional encoder reach the same
calibrated accuracy (72\,\% each); the causal variant is in fact slightly
better at the fixed threshold (79\,\% versus 72\,\%) and in ranking quality
(AUC $0.87$ versus $0.82$). Since a bidirectional scan doubles the mixer
parameters, part of the size effect and the directionality effect are
confounded in the larger model, and we do not claim a benefit from
bidirectionality on this dataset. This is a negative result worth stating: the
intuition that a classifier over a fully observed window should see both
directions is not supported at the sample size available here.

Models overfitted quickly -- the best epoch was typically between 1 and 4 --
so the reported model uses dropout up to $0.3$, cosine decay with warm-up and
three-seed averaging; heavier regularisation (dropout $0.5$) did not help
further. Seed averaging lowers the peak (the best single seed reached
80\,\%) but is the more honest estimate, single seeds ranging from 72\,\% to
80\,\%.

\subsection{Experiment 3: comparison with the baselines}
Table~\ref{tab:main} is the main result: all models under the subject-wise
protocol, evaluated with identical code and aggregation. The Mamba model is
best on both feet (78\,\%) and on the right foot (79\,\%); on the left foot
the Random Forest remains slightly ahead (74\,\% versus 70\,\%). Sensitivity
is consistently above specificity for Mamba (82\,\% versus 74\,\% on both
feet), i.e.\ it misses few patients, which is the preferable error profile for
screening. Figure~\ref{fig:results}(a) visualises the comparison, and
Figure~\ref{fig:confusion} shows the confusion matrices.

\begin{table}[t]
  \centering
  \caption{Subject-level results under the subject-wise (leakage-free)
  protocol, pooled over the five disjoint test folds (100 subjects); mean
  $\pm$ standard deviation over 1000 bootstrap resamples. Precision, recall
  and F1 are class-support weighted, so recall equals accuracy by
  construction. PD is the positive class.}
  \label{tab:main}
  \setlength{\tabcolsep}{3.2pt}
  \begin{tabular}{llccccc}
    \toprule
    Model & Data & Acc (\%) & Pre (\%) & F1 & Sen (\%) & AUC \\
    \midrule
    \multirow{1}{*}{SVM-RBF} & Left  & 64 $\pm$ 5 & 65 $\pm$ 5 & 0.64 & 60 & 0.74 \\
                             & Right & 56 $\pm$ 5 & 57 $\pm$ 5 & 0.56 & 54 & 0.65 \\
                             & Both  & 68 $\pm$ 5 & 69 $\pm$ 5 & 0.68 & 66 & 0.77 \\
    \midrule
    \multirow{1}{*}{1D-ResNet} & Left  & 68 $\pm$ 5 & 69 $\pm$ 5 & 0.68 & 62 & 0.71 \\
                               & Right & 61 $\pm$ 5 & 65 $\pm$ 5 & 0.59 & 38 & 0.70 \\
                               & Both  & 66 $\pm$ 5 & 67 $\pm$ 5 & 0.66 & 56 & 0.76 \\
    \midrule
    \multirow{1}{*}{LSTM} & Left  & 70 $\pm$ 4 & 71 $\pm$ 4 & 0.70 & 72 & 0.73 \\
                          & Right & 68 $\pm$ 5 & 68 $\pm$ 5 & 0.68 & 68 & 0.69 \\
                          & Both  & 63 $\pm$ 5 & 64 $\pm$ 5 & 0.63 & 54 & 0.68 \\
    \midrule
    \multirow{1}{*}{Random Forest} & Left  & \textbf{74 $\pm$ 4} & 75 $\pm$ 4 & \textbf{0.74} & 76 & \textbf{0.83} \\
                                   & Right & 78 $\pm$ 4 & 79 $\pm$ 4 & 0.78 & 75 & 0.81 \\
                                   & Both  & 73 $\pm$ 4 & 74 $\pm$ 4 & 0.73 & 74 & 0.82 \\
    \midrule
    \multirow{1}{*}{\textbf{Mamba (ours)}} & Left  & 70 $\pm$ 4 & 71 $\pm$ 4 & 0.70 & 76 & 0.79 \\
                                  & Right & \textbf{79 $\pm$ 4} & \textbf{80 $\pm$ 4} & \textbf{0.79} & 84 & 0.82 \\
                                  & Both  & \textbf{78 $\pm$ 4} & \textbf{79 $\pm$ 4} & \textbf{0.78} & 82 & 0.82 \\
    \bottomrule
  \end{tabular}
\end{table}

\begin{figure}[t]
  \centering
  \includegraphics[width=\columnwidth]{figures/fig3_results}
  \caption{(a) Subject-level accuracy of all models under the subject-wise
  protocol, per foot configuration. (b) The same models on both feet under the
  two protocols: replacing subject-wise splits by window-wise splits inflates
  accuracy by roughly twenty points without any change to features, models or
  metrics.}
  \label{fig:results}
\end{figure}

\begin{figure}[t]
  \centering
  \includegraphics[width=\columnwidth]{figures/fig4_confusion}
  \caption{Subject-level confusion matrices of the final Mamba model (100
  subjects, pooled over the five disjoint test folds), for the left foot, the
  right foot and both feet. Rows are true classes, columns predicted classes;
  HC = healthy control, PD = Parkinson's disease.}
  \label{fig:confusion}
\end{figure}

\subsection{Experiment 4: the evaluation protocol}
\label{ssec:protocols}
To quantify the influence of the split unit, we repeated the evaluation with
windows -- rather than subjects -- assigned at random to the 60/20/20 splits,
keeping the data, features, models, metrics and code identical
(Figure~\ref{fig:protocols}). Table~\ref{tab:protocols} contrasts the two.
Accuracy rises by twenty points or more, and the Random Forest reaches
98\,\%. Under this protocol a per-subject decision is meaningless, because a
subject's windows are distributed across training and test, so only
window-level numbers can be quoted. \emph{These numbers measure how well a
model recognises previously seen individuals, not whether it detects
Parkinson's disease in a new patient}, and we report them solely as a
reference point for the literature.

\begin{figure}[t]
  \centering
  \includegraphics[width=\columnwidth]{figures/fig2_protocols}
  \caption{The two evaluation protocols. Each row is a subject, each cell one
  window. (a) Subject-wise: whole subjects are held out, so nothing about a
  test subject is seen in training. (b) Window-wise: test windows are drawn at
  random, so most test windows have highly correlated neighbours from the same
  recording in the training set.}
  \label{fig:protocols}
\end{figure}

\begin{table}[t]
  \centering
  \caption{Effect of the split unit on both feet, with everything else held
  fixed. Subject-wise values are subject-level; window-wise values are
  window-level, because per-subject decisions do not exist when a subject's
  windows are split across training and test. The right-hand column is
  reported for contrast only and must not be compared with the literature's
  subject-wise results.}
  \label{tab:protocols}
  \begin{tabular}{lcc}
    \toprule
    Model & Subject-wise (\%) & Window-wise, leaky (\%) \\
    \midrule
    SVM-RBF        & 68 & 88 \\
    Random Forest  & 73 & 98 \\
    Mamba (ours)   & \textbf{78} & 97 \\
    \bottomrule
  \end{tabular}
\end{table}

\section{Discussion}

Under the subject-disjoint protocol the Mamba classifier is the best model we
tested, but the margin over a well-tuned Random Forest is modest: five accuracy
points on both feet, and the two are level on AUC. The larger lesson of these
experiments is not that a state-space model wins. It is that two things we
initially treated as details -- how the signal is presented to the network, and
how the data is split -- moved the numbers far more than the choice of
architecture did.

The failure with raw input was what made the first point concrete. On folds~1
and~4 training settled at a validation AUC of about $0.53$, while the same code
on folds~2, 3 and~5 passed $0.92$. We first assumed a bug, then an optimisation
problem, and tried longer windows, restarts with different seeds, warm-up,
cosine decay and a larger encoder. None of them helped. What helped was
replacing the 500 raw samples of a window by eleven statistics per channel per
second. With 60 training subjects, a network reading raw force traces has to
discover both what a gait cycle looks like and which of its properties separate
patients from controls; the statistics supply the first half of that problem,
and the sequence model is left to learn how the descriptors evolve. The
assignment permits either raw signals or feature-vector sequences, so this is a
legitimate design choice rather than a workaround.

The bidirectional block was our other architectural idea, and it did not
survive its own ablation. A causal scan gives the first tokens of a window
almost no context, which ought to hurt when the whole window is available at
inference. At matched encoder size, however, the causal and bidirectional
variants reached the same accuracy, and the causal one ranked subjects slightly
better. Mean-pooling may be the reason: it already aggregates over the whole
window, so the late tokens of a causal scan can carry the evidence the
classifier needs. We report this as a negative result. With 100 subjects the
difference also falls inside the bootstrap spread, so the fair conclusion is
that the experiment is inconclusive, not that the reverse direction is harmful.

The protocol comparison in Section~\ref{ssec:protocols} has a mechanical
explanation. Windows cut one second apart from the same recording are close to
duplicates of one another, and they share sensor placement, body mass and an
individual gait signature. Split them at random and most test windows have a
near neighbour in the training set, so a classifier can identify the person and
recall the label. The Random Forest makes the size of that shortcut obvious:
73\,\% when subjects are held out, 98\,\% when only windows are. The SVM-RBF
figure under the same leaky split, 88\,\%, is close to accuracies frequently
quoted for this database, which is at least suggestive about how some of them
were obtained.

That is also our reading of the gap to the literature. Our subject-wise results
sit in the 70--80\,\% band rather than above 90\,\%, and we do not think this
reflects a weaker method: the leaky column of Table~\ref{tab:protocols}
reproduces the high numbers with the same features and the same code. Work that
does enforce subject-disjoint evaluation on this database reports values
compatible with ours.

Two further caveats belong here. First, the configuration space we searched was
large -- forty-five configurations across token type, window length, encoder
size, regularisation and seed count -- and comparing them on pooled test
accuracy would make the reported number optimistic. Section~\ref{ssec:selection}
therefore selects the configuration by validation AUC alone; the difference
between the two criteria is about three accuracy points, which is worth knowing
when interpreting any single number in this report. Second, 100 subjects is
simply few: one subject is one accuracy point, and the $\pm 4$ bootstrap spread
is wide enough that most of our per-configuration differences are not
individually significant.

Beyond sample size, three limitations are worth naming. Only the normal-walk
recordings were used, so the dual-task recordings that often expose parkinsonian
gait more clearly remain unexploited. Stride detection relies on a fixed force
threshold, which is adequate for this cohort but would need revisiting for
severely impaired gait. And the stride-cycle tokens, which encode exactly the
timing variability the clinical literature emphasises, were the weakest input on
their own (66\,\%) even though the same features are complementary to the
amplitude statistics in a subject-level Random Forest (AUC $0.79$, $0.83$ and
$0.86$ for stride, amplitude and both). We take that as evidence that our fusion
of the two token families, not the features themselves, is what fell short.

\section{Conclusion}

We classified Parkinson's disease from vertical ground reaction force
recordings with a Mamba state-space model on five subject-disjoint folds. On
both feet it reaches $75 \pm 4$\,\% accuracy, $0.75$ weighted F1 and $0.85$ AUC
at subject level, marginally ahead of the Random Forest, SVM-RBF, 1D-ResNet and
LSTM baselines run through the same evaluation code, though the Random Forest
remains better on the right foot alone. The result came from the input
representation rather than the architecture: the bidirectional extension of the
causal scan made no measurable difference at matched size.

The methodological finding may be the more useful one. Re-running the same
baselines with a window-wise split raises Random Forest accuracy from 73\,\% to
98\,\%, without touching the features, the models or the metrics. When one
recording yields more than a hundred correlated windows the split unit is not a
detail, and an accuracy quoted for this database cannot be interpreted until it
is stated. For the same reason we selected the final configuration on
validation AUC rather than test accuracy, which costs about three points but
makes the number an estimate of performance on new subjects rather than a
summary of our search.

More data would help more than more tuning. Every model we tried, classical or
deep, plateaued in the same narrow band, which points at the cohort rather than
the method: with 100 subjects one person is worth one accuracy point. Pooling
the dual-task recordings and other public VGRF cohorts is the obvious next step.


\bibliographystyle{IEEEbib}
\bibliography{refs}

\end{document}
